\section{The Saints and the Holy Angels}\label{sec:saints}

The saints have lived what the Fathers and Doctors taught. From the
thirteenth century to the nineteenth the Church's holy men and women kept
company with the angels as fellow citizens of the city of God: they fasted
in honour of Saint Michael, they spoke with their guardian angels as with
friends, they were wounded and purified by the fire of the Seraphim, and
they learned to tell the good angel from the enemy who disguises himself as
an angel of light. Their witness runs along four lines. The first is
Michael, prince of the heavenly host and bearer of souls to God, whom
Francis honoured with a Lent of forty days, whom Joan of Arc saw with her
bodily eyes, and whom Alphonsus sends to the bedside of the dying. The
second is the guardian angel, given to each believer and made visible to a
few, as to Frances of Rome and Gemma Galgani, and sung by Thérèse of the
Child Jesus and by Newman. The third is the burning love of the highest
orders, the Seraph of La Verna and the angel of Teresa's transverberation,
in which the saints found the pattern of their own union with God. The
fourth is the angelic life itself, the purity by which Thomas Aquinas became
the Angelic Doctor and Gertrude and Mechtild were drawn into the choirs.

Bernard of Clairvaux, who stands at the head of this company, is treated
with the medieval Doctors (\S\ref{sec:fathers-after}); Aquinas's own
teaching on the guardians is given in its place in the treatise
(\S\ref{sec:q113}), and the Church's rule for the honour of the angels in
\S\ref{sec:devotion}.

\sectionguard\subsection{Francis of Assisi: The Lent of Saint Michael and
the Seraph}\label{sec:sa-francis}

Bonaventure, seventh successor of Francis in the government of the Friars
Minor, wrote the life through which the Church has read the Poverello ever
since. Its prologue sets Francis at once among the angels. He preached
the gospel of peace, ``himself an Angel of the true peace''; he was
``appointed unto angelic ministries; thereafter, wholly set on fire by the
kindling of the Seraph, and, like the prophet, borne aloft in a chariot of
fire'' (\emph{Legenda maior}, prol.~1, trans.\ E. Gurney Salter, 1904).
Bonaventure then applies to him the angel of the sixth seal: ``he is
thought to be not unmeetly set forth in the true prophecy of that other
friend of the Bridegroom, the Apostle and Evangelist John, under the
similitude of the Angel ascending from the sunrising and bearing the seal of
the Living God'' (prol.~1; Apoc 7:2). The reason is his life: ``living among
men, he was an imitator of the purity of the Angels'' (prol.~2).
Bonaventure presents the identification as a belief to be ``faithfully and
devoutly held'' (prol.~2).

The angels frame the whole life. When Francis began to repair the little
church of Saint Mary, he chose to live there because heaven had already
chosen it: ``Perceiving that Angels ofttimes visited it,---according unto the
name of that church, that from old time was called Saint Mary of the
Angels,---he abode there by reason of his reverence for the Angels, and his
especial love for the Mother of Christ'' (II.8). There he began, there he
grew, there he died, and dying ``commended it unto the Brethren as a place
most beloved of the Virgin.'' Of his prayer Bonaventure
writes that Francis had learned, ``like the heavenly spirits on Jacob's
ladder, to be ever ascending toward God, or stooping toward his neighbour''
(XIII.1; Gen 28:12), dividing his time between the labour of preaching and
``the peaceful ecstasies of contemplation.''

\subsubsection{Two Lents for the angels}

Francis kept several Lents in the year, and two of them belonged to the
angels. Bonaventure gives the order of his devotions in the ninth chapter,
on his burning love. After Christ he trusted most in the Mother of God, and
fasted in her honour from the feast of Peter and Paul to the Assumption.
Then:
\begin{quote}
He was bound by ties of inseparable affection unto the Angelic spirits that
do glow with wondrous fire to approach God, and in the kindling of elect
souls, and out of devotion unto them he would fast for forty days from the
Assumption of the glorious Virgin, remaining instant in prayer throughout
that time. Unto the Blessed Michael Archangel,---inasmuch as his is the
ministry of bringing souls before God,---he cherished an especial love and
devotion, by reason of the ardent zeal that he had for the salvation of all
such as should be saved. (IX.3)
\end{quote}
The Latin gives the reason for the love of Michael in a single clause:
\latin{Beato autem Michaeli Archangelo, eo quod animarum repraesentandarum
haberet officium, speciali erat amore devotior} --- to blessed Michael the
Archangel, because his was the office of presenting souls, he was devoted
with a special love (\emph{Legenda S.\ Francisci} IX.3, Quaracchi,
\emph{Opera} VIII). Francis loved the angels because they burn, and he loved
Michael because Michael brings the saved home; the zeal of the preacher for
souls and the office of the archangel who presents them are one desire. The
forty days from the Assumption end at Michaelmas, and the Quaracchi editors
note that the thirteenth chapter calls the same fast the Lent in honour of
Saint Michael. Beside Bonaventure's clause they set the responsory after the
fifth lesson of the Roman Breviary for 29 September, \latin{Cui tradidit Deus animas Sanctorum, ut
perducat eas in paradisum exsultationis}, and the third antiphon at Lauds,
\latin{Archangele Michael, constitui te principem super omnes animas
suscipiendas} (IX.3, n.~10 in the Quaracchi edition): the office of Michael
that Francis honoured is the office the Church sings.

\subsubsection{The Seraph of La Verna}

The Lent of Saint Michael was the setting of the greatest grace of
Francis's life. Two years before his death he went up the mountain of La
Verna, and ``When, according unto his wont he began to keep a Lent there,
fasting, in honour of Saint Michael Archangel, he was filled unto
overflowing, and as never before, with the sweetness of heavenly
contemplation'' (XIII.1). The Latin names the fast exactly:
\latin{Quadragesimam ibidem ad honorem sancti Archangeli Michaelis ieiunare
coepisset}. Bonaventure notes that even the birds of the mountain seemed to
welcome him, and a falcon woke him for the night office; ``there would seem
to have been a divine omen \dots\ inasmuch as that praiser and worshipper of
God, upborne on the wings of contemplation, was at that very place and time
to be exalted by the vision of the Seraph'' (VIII.10).

Three times the book of the Gospels, opened in the name of the Trinity, fell
open at the Passion, and Francis understood that he must be made like
Christ in suffering before he died (XIII.2). Then came the vision:
\begin{quote}
When, therefore, by seraphic glow of longing he had been uplifted toward
God, and by his sweet compassion had been transformed into the likeness of
Him Who of His exceeding love endured to be crucified,---on a certain
morning about the Feast of the Exaltation of Holy Cross, while he was
praying on the side of the mountain, he beheld a Seraph having six wings,
flaming and resplendent, coming down from the heights of heaven. When in
his flight most swift he had reached the space of air nigh the man of God,
there appeared betwixt the wings the Figure of a Man crucified, having his
hands and feet stretched forth in the shape of a Cross, and fastened unto a
Cross. Two wings were raised above His head, twain were spread forth to
fly, while twain hid His whole body. (XIII.3)
\end{quote}
The Seraph is the one Isaias saw, whose six wings covered face and feet and
bore him in flight (Isa 6:2); the Latin says the wings were \latin{tam
ignitas quam splendidas}, as fiery as they were bright. Francis rejoiced and
grieved at once: he rejoiced ``at the gracious aspect wherewith he saw
Christ, under the guise of the Seraph, regard him,'' and the crucifixion
pierced him with pity; he marvelled, ``knowing that the infirmity of the
Passion doth in no wise accord with the immortality of a Seraphic spirit.''
The Lord gave him the meaning: he was to be ``wholly transformed into the
likeness of Christ Crucified, not by martyrdom of body, but by enkindling of
heart'' (XIII.3). As the vision passed it left a fire in his heart and the
marks of the nails in his hands and feet and the wound in his side.

Bonaventure keeps a silence where Francis kept one. The saint told the
vision only when Brother Illuminato persuaded him that such secrets were
shown him for the sake of others also, ``adding that He who had
appeared unto him had said some words the which, so long as he lived, he
would never reveal unto any man. Verily we must believe that those
utterances of that holy Seraph marvellously appearing on the Cross were so
secret that perchance it was not lawful for a man to utter them'' (XIII.4;
2 Cor 12:4). And the fast ended on the angels' day: ``when the Feast of
Saint Michael Archangel came, this angelic man, Francis, descended from the
mountain, bearing with him the likeness of the Crucified, engraven, not on
tables of stone or of wood, by the craftsman's hand, but written on his
members of flesh by the finger of the Living God'' (XIII.5). In
Bonaventure's telling the vision belongs to the angels by its form and to
Christ by its content: the \latin{seraphicis desideriorum ardoribus}, the
seraphic ardours of desire that lifted Francis toward God (XIII.3), were
answered by Christ under the likeness of the order that burns, and the love
of Christ crucified was the fire that marked his servant. The fire of La Verna marks the whole Franciscan tradition, and
Bonaventure's own works are printed under the name of the Seraphic Doctor
(\latin{Doctoris seraphici S.\ Bonaventurae opera omnia}).

\sectionguard\subsection{Bonaventure: The Six Wings and the Sermons on the
Angels}\label{sec:sa-bonaventure}

\subsubsection{The Seraph as the way of ascent}

Thirty-three years after the death of Francis, Bonaventure, then Minister
General, went up La Verna himself. He tells what happened in the prologue of
the \emph{Journey of the Mind into God}. He had gone there \latin{amore
quaerendi pacem spiritus}, for love of seeking peace of spirit, and while he
turned over in his mind certain ascents of the mind into God, \latin{inter
alia occurrit illud miraculum, quod in praedicto loco contigit ipsi beato
Francisco, de visione scilicet Seraph alati ad instar Crucifixi} --- among
other things there came to him that miracle which befell blessed Francis in
that same place, the vision of the winged Seraph in the likeness of the
Crucified. At once it seemed to him that the vision showed both the
father's rapture in contemplation and the way to it (\emph{Itinerarium},
prol.~2, Quaracchi, \emph{Opera} V):
\begin{quote}
\latin{Nam per senas alas illas recte intelligi possunt sex illuminationum
suspensiones, quibus anima quasi quibusdam gradibus vel itineribus
disponitur, ut transeat ad pacem per ecstaticos excessus sapientiae
christianae. Via autem non est nisi per ardentissimum amorem Crucifixi.}
(prol.~3)
\end{quote}
By those six wings may rightly be understood six uplifting illuminations by
which the soul is disposed, as by steps or stages, to pass into peace
through the ecstatic transports of Christian wisdom; and the way is nothing
but the most burning love of the Crucified. The same love so absorbed the
mind of Francis \latin{quod mens in carne patuit}, that his mind showed in
his flesh when he bore the stigmata for two years before his death. The six
wings of the Seraph therefore signify six illuminations \latin{quae a
creaturis incipiunt et perducunt usque ad Deum, ad quem nemo intrat recte
nisi per Crucifixum} --- which begin from creatures and lead all the way to
God, whom no one rightly enters except through the Crucified (prol.~3). The
ascent of the whole \emph{Itinerarium} is laid out on the Seraph of La
Verna: \latin{Effigies igitur sex alarum seraphicarum insinuat sex
illuminationes scalares}, the figure of the six seraphic wings signifies six
illuminations that are the rungs of a ladder.

\subsubsection{The ladder of the hierarchies}

Bonaventure's sermons for the feast of Saint Michael preach the angelic
hierarchy to the brethren. The Quaracchi editors print five among his
sermons for the saints
(\emph{Sermones de sanctis}, \latin{De sanctis Angelis} 1--5, Quaracchi,
\emph{Opera} IX, 609 ff.); the first is edited from a thirteenth-century
Troyes manuscript whose many faults the editors corrected as best they
could, but its design is plain. He takes Jacob's ladder for his text (Gen
28:12) and begins with a prayer: \latin{Habemus loqui de hierarchiis
angelicis, quae sunt in caelo supra sensum et rationem; ideo in principio
rogemus Dominum, quod det mihi Spiritum suum} --- we have to speak of the
angelic hierarchies, which are in heaven above sense and reason; therefore
at the beginning let us ask the Lord to give me his Spirit. The ladder is
one order spanning heaven and earth:
\begin{quote}
\latin{Scala ista, cuius pars est in caelo et pars in terra, significat
ordinationem angelicae hierarchiae, cuius pars est in angelicis spiritibus et
Sanctis et pars in hominibus et animabus sanctis.} (s.~1, p.~610)
\end{quote}
This ladder, part in heaven and part on earth, signifies the ordering of the
angelic hierarchy, part of which is in the angelic spirits and the saints
and part in men and holy souls. The heavenly part has three hierarchies of
three orders each, which Bonaventure takes from Dionysius, and the earthly
part answers it with three hierarchies of the Church: patriarchs, prophets,
and apostles in the primitive Church; martyrs, confessors, and virgins in
the Church advancing; prelates, contemplatives, and actives in the Church
spread abroad.

He then reads the visible heaven as a mirror of the invisible. The first
hierarchy answers to the nobility of the heavens, the second to their
power, the third to their abundant influence on the earth, and the whole to
their ordered harmony. His summary gives each order its office and binds
the first three to the Persons of the Trinity:
\begin{quote}
\latin{In prima hierarchia sunt Throni, qui respondent potestati Patris, et
Cherubim, qui respondent Filio, et seraphicus ordo, qui respondet suavitati
Spiritus sancti. In secunda hierarchia sunt tres ordines, scilicet ordo
Dominationum, Virtutum et Potestatum; et Dominationum est imperare;
Virtutum, exsequi; et Potestatum, triumphare.} (s.~1, pp.~612--613)
\end{quote}
The Thrones answer to the power of the Father, the Cherubim to the Son, the
seraphic order to the sweetness of the Holy Spirit; of the second hierarchy
the Dominations command, the Virtues execute, the Powers triumph. The third
hierarchy, which looks to the Church on earth and flows into it, feeds men
with a threefold bread: the Principalities lead us by the hand, the
Archangels teach us the secrets of God, and the Angels support and guard us
in prosperity and adversity. Above the nine orders stand two more: \latin{Et
supra se hierarchia angelica habet Virginem gloriosam; et beata Virgo supra
se habet Iesum Christum, Filium suum} --- the angelic hierarchy has above it
the glorious Virgin, and the blessed Virgin has above her Jesus Christ, her
Son (p.~612). Their song is one music: \latin{Concentus caeli est harmonia
laudis angelicae, quia, sicut in cithara sonat quaelibet chorda secundum
proportionem et virtutem suam, sic secundum dona (et) virtutes in angelica
hierarchia novem ordinum est exsultationum, iucunditatum et laudum mira
consonantia} --- the harmony of heaven is the concord of angelic praise;
as on a harp each string sounds according to its measure and strength, so in
the hierarchy of nine orders there is, according to their gifts and powers,
a wonderful consonance of exultation, gladness, and praise.

The sermon ends with an exhortation. Michael, Bonaventure says, was once set
over the synagogue and is now set over the Church, \latin{ideo eius officium
celebratur per universam Ecclesiam} --- therefore his feast is kept through
the whole Church (p.~613; Dan 12:1). The angels offer everything, and we
refuse it: \latin{Angeli dicunt nobis: Suscipite exempla, documenta et
servitia nostra et ministeria; et nos sumus rusticani, quia claudimus eis
ostium} --- the angels say to us, take our examples, our teachings, our
services and ministries; and we are boors, for we shut the door on them
(p.~612). \latin{Magnum est confidere in Angelis}, it is a great thing to
trust in the angels; and he turns on the learned who speak of heavenly
things without love: \latin{Loquuntur de caelestibus et affectum habent in
terrenis}, they talk of heavenly things and their affection is on earthly
ones.

\subsubsection{Remedies, examples, and benefits}

The fourth sermon takes the words of Christ, ``their angels in heaven always
see the face of my Father'' (Matt 18:10), and divides the angels' office in
two: to minister to men and to contemplate God. \latin{Angeli dicuntur
nostri propter tria, quae ministrant nobis, scilicet remedia, exempla et
beneficia} --- they are called our angels for three things they minister to
us, remedies, examples, and benefits (s.~4, p.~620). The three remedies
answer the three wounds of sin, proneness to evil, difficulty in good, and
distance from God, and each is ministered by an archangel named in
Scripture. \latin{Primum remedium ministrat Raphael, qui interpretatur
medicina Dei}: Raphael, the medicine of God, healed Tobias of blindness and
Sara of the demon, that is, our ignorance and our concupiscence (Tob 6, 11).
\latin{Secundum remedium ministrat Gabriel, qui interpretatur fortitudo Dei,
quia facit fortes contra difficultatem}: Gabriel, the strength of God,
announced the birth of the Forerunner under the law of marriage and the
birth of Christ from a Virgin, and so strengthens us to keep both the
precepts and the counsels. \latin{Tertium remedium ministrat Michael, qui
interpretatur Quis ut Deus. Nam sua ministratio praesentes animas Deo
facit}: Michael, \emph{Who is like God}, makes souls present to God, as the
Church sings of him, that he \latin{introducit animas in paradisum
exsultationis}. His name itself is the cure of our exile: \latin{quis est
tam desiderandus ut Deus?} --- who is so much to be desired as God? And so,
Bonaventure concludes, \latin{tres Angelos solum ex nomine designat sacra
Scriptura, scilicet Raphaelem, Gabrielem et Michaelem} --- Holy Scripture
names only three angels, Raphael, Gabriel, and Michael (s.~4, pp.~620--621),
the rule the Church keeps for the angels' names (\S\ref{sec:devotion}).

The four examples are reverence toward God, when the angels fall on their
faces before the throne (Apoc 7:11), purity in themselves, concord among
themselves, and clemency toward men, since they carried the beggar Lazarus
to Abraham's bosom and rejoice over one sinner doing penance (Luke 16:22;
15:10). For their purity Bonaventure gives the definition he took from
Dionysius, the angel as \latin{imago Dei manifestativa luminis occulti,
speculum purum}, an image of God that manifests the hidden light, a pure
mirror, and adds the saying he read under Jerome's name, \latin{Semper est
Angelis cognata virginitas}, virginity is ever kin to the angels.
The five benefits are cleansing, as the Seraph cleansed Isaias's lips with
the coal (Isa 6:6--7), enlightenment, inflaming with love, protection from
the demons (Ps 33[34]:8; Acts 12:11), and the general custody of souls:
\latin{nam generales custodes sunt animarum Angeli} (p.~621; Ps 90[91]:11).
The fifth sermon, on that last verse, carries the doctrine into the
Trinity: \latin{hierarchia nihil aliud est quam sacer principatus}, a
hierarchy is nothing but a sacred rule, and there are three, the
supercelestial in the divine Persons, the celestial in the angels, and the
subcelestial in men, the angels standing between as the aeon stands between
eternity and time (s.~5).

\sectionguard\subsection{Thomas Aquinas, the Angelic Doctor}\label{sec:sa-aquinas}

The Church calls the author of the treatise on the angels by their name, and
his early biographers tie the name to an angelic grace in his youth. When
Thomas took the Dominican habit against his family's will, his brothers held
him prisoner and tried at last to overcome him with a woman sent into his
room. William of Tocco tells what followed. Thomas drove her out with a
brand from the fire, traced a cross on the wall with it, and prostrate on the ground begged God with
tears for the girdle of perpetual virginity; and as he prayed he fell asleep:
\begin{quote}
\latin{ecce ad eum duo Angeli coelitus missi sunt, qui asserentes eum a
Domino exauditum et de pugna tam difficili obtinuisse triumphum,
stringentesque ipsum hinc inde in renibus, dixerunt: Ecce ex parte Dei te
cingimus, quod petivisti, cingulo castitatis, quod nulla possit de cetero
impugnatione dissolvi, et hoc quod humanae virtutis haberi non potest ex
merito, tibi conceditur divinae largitatis ex dono.} (Tocco, \emph{Vita S.\
Thomae} c.~10, ed.\ Prümmer, \emph{Fontes vitae S.\ Thomae Aquinatis})
\end{quote}
Behold, two angels were sent to him from heaven, who assured him that the
Lord had heard him and that he had won the victory in so hard a fight, and
binding him about the loins on either side they said: behold, on God's part
we gird you, as you asked, with the girdle of chastity, which no assault
shall ever loose; and what cannot be had by the merit of human virtue is
granted you by the gift of divine bounty. The pressure of the angelic touch
was so painful that he woke with a cry; he kept the grace secret and told
it only to his companion at the end of his life. The life by Peter
Calo gives the angels' words as \latin{Cingimus te cingulo perpetue
virginitatis, quod nullatenus de cetero dissolvetur} and records that at the
hour of his death Thomas told Brother Reginald \latin{quod ex tunc non sensit
nec minimum carnis motum}, that from that time he had not felt the least
movement of the flesh (Calo, \emph{Vita} 7, ed.\ Prümmer, p.~24). Prümmer
notes that Tocco attested the same story on oath in the process of
canonization.

Tocco closes the chapter with a cry of praise in which the girding and the
title meet. Blessed the youth, he writes, \latin{cui Angelica societas, dum
castitate cingitur, non negatur, qui meruit fieri puritate Angelicus} ---
to whom the company of angels is not denied, since he is girded with
chastity, and who merited to become angelic by purity. Pius XI made the
same connection in the encyclical \emph{Studiorum Ducem} (29 June 1923) for
the sixth centenary of the canonization. Among the virtues of Thomas,
\latin{prima omnium occurrit ea virtus, unde quaedam cum angelicis naturis
visa est esse Thomae similitudo; castimoniam dicimus}, the first that meets
us is the one by which Thomas seemed to bear a likeness to the angelic
natures, namely chastity, which he kept unharmed in the gravest danger, so
that \latin{dignus est habitus quem mystica zona angeli cingerent} --- he was
held worthy to be girded by angels with a mystical belt. The Pope recalls
that Pius V inscribed him among the Doctors \latin{rato Angelici titulo},
with the title of Angelic confirmed, and himself adds that he should be
called \latin{non modo Angelicum, sed etiam Communem seu universalem Ecclesiae
Doctorem}. The purity and the doctrine are one gift: \latin{si Thomae
pudicitia \dots\ cecidisset, verisimile est nequaquam Ecclesiam suum
Doctorem Angelicum habituram fuisse} --- had the purity of Thomas fallen,
the Church would likely never have had her Angelic Doctor.

The grace became an institution. Pius XI commended to the bishops the
\latin{Militiae Angelicae societatem}, the Angelic Warfare founded to guard
chastity under Thomas's protection, confirmed the indulgences his
predecessors had granted it, and allowed its members to wear, instead of the
girdle, a medal bearing on its face the image \latin{sancti Thomae cum Angelis
ei zonam accingentibus}, of Saint Thomas with the angels girding him. At the
end of the letter, so that studies under Thomas's guidance might bear more
fruit, he appended the form of prayer \latin{qua ipse utebatur}, which
Thomas himself used; it begins with the angels:
\latin{Creator ineffabilis, qui de thesauris sapientiae tuae tres Angelorum
hierarchias designasti, et eas super caelum empyreum miro ordine collocasti}.
The Doctor who set out the three hierarchies in the \emph{Summa}
(\S\ref{sec:hierarchies}) addressed the Creator in prayer as the one who had
designed them from the treasures of his wisdom and set them above the
empyrean heaven in wonderful order, and asked of that same Source of light a
ray of its clarity on the darkness of his mind.

\sectionguard\subsection{Gertrude and Mechtild}\label{sec:sa-helfta}

Two Benedictine nuns, Gertrude, called the Great, and Mechtild, whose
prayers Gertrude sought and to whom, her book records, Gertrude's sanctity
was revealed, received revelations in which the choir on earth and the
choirs of heaven sing together. Gertrude's \emph{Insinuationes divinae
pietatis} were compiled from her papers by a religious of her monastery, all
but the second book, which she wrote herself; the Kenmare translation of
1870 renders the second to fifth books as its second to fifth parts.
Mechtild's revelations were rendered into English in a London selection of
1875, taken from the five books of her \emph{Spiritual Grace}.

\subsubsection{Gertrude and the feast of Saint Michael}

The angels in Gertrude's revelations are liturgical. At Mass one day, ``it
appeared to her that her guardian angel took her in his arms as if she were
a little child, at the Kyrie Eleison, and presented her to God the Father,
to receive His benediction, saying: `Eternal Father, bless Thy little
child'\,'' (\emph{Life and Revelations}, pt~III ch.~21). The Father was
silent, Gertrude sank into the knowledge of her nothingness, and the Son
clothed her with the merits of his life, so that the angel's request was
granted through Christ. At the Christmas chapter the guardian of each nun
presented her heart to God, and the choirs came according to the virtue of
each: ``When those who loved God with the greatest fervour offered Him their
hearts, the angels of the choir of Seraphim, who attended our Lord and
supported Him, disposed those religious for their offering. When the hearts
of those who were most enlightened in the knowledge of God were offered, the
angels of the choir of Cherubim came to present their homage'' (pt~IV
ch.~2). At Vespers in the same season she saw angels flying round the
convent singing the verse of the hymn with the nuns, and learned ``that when
angels are present at our solemnities, they pray to God that those who
imitate them in their devotion may imitate them in purity of body and
soul'' (pt~IV ch.~2); the psalmist's word, ``I will sing praise to thee in
the sight of the angels'' (Ps 137[138]:1), was fulfilled in her choir.

Her chapter for Michaelmas is a lesson in the honour owed to the angels. As
the feast approached, Gertrude prepared for Communion ``by meditating on the
care which the angels had of her, by the Divine command, notwithstanding her
unworthiness,'' and wishing to make some return she offered the Sacrament
itself in their honour: ``I offer Thee this most august Sacrament, O most
loving Lord, for Thy eternal glory, in honour of the princes of Thy kingdom,
and for the increase of their felicity and beatitude'' (pt~IV ch.~54). The
honour of the angels passes through the Eucharist to God, and returns to the
soul. Each choir bowed before her and promised its own help:
\begin{quote}
``Thou hast indeed honoured us by this oblation, and we will therefore
guard thee with special care;'' the guardian angels adding: We will guard
thee night and day with ineffable joy, and will prepare thee for thy Spouse
with the utmost vigilance. (pt~IV ch.~54)
\end{quote}
She recognized her own guardian, who ``appeared to her as a prince
magnificently attired, standing between her soul and God, and endeavouring
to unite and elevate her soul to God.'' The Archangels promised: ``we will
discover to you the Divine secrets, according to your capacity of receiving
them''; the Virtues, ``We will assist you in your labours, writings, and
meditations for the glory of God''; the Dominations, to offer for her the
honour she owed to God's sovereignty; the Principalities, to present her to
the King adorned according to his heart; the Powers, to ``continually remove
every impediment, whether exterior or interior, which might interrupt His
Divine communications.'' The Powers gave the reason: ``For the prayers of
one loving soul prevail more with God, both for the living and the dead,
than the prayers of a thousand souls who love less'' (pt~IV ch.~54). Each
order does for her what its name says it does for the Church, the
Archangels revealing, the Virtues working, the Powers restraining.

\subsubsection{Mechtild's golden staircase}

Mechtild saw the nine orders as a ladder that men climb. ``The handmaid of
Christ beheld a golden staircase composed of nine steps, with a multitude of
angels surrounding it on either side, so that on the first step were the
Angels, and on the second the Archangels, and so upwards, in such a way that
on each step there was placed an angelic order. And she understood from
heaven, that by staircase is signified the conversation of men'' (\emph{Select
Revelations}, ch.~16, ``Of the Angels, and how men are made their
companions''). Each step joins a way of life to an order. Whoever
``faithfully, and humbly, and devotedly ministereth in the Church of God, and
also for God's sake serveth the infirm, or pilgrims, or the poor'' stands on
the first step and is made equal with the Angels; those who wait on God in
prayer serve with the Archangels; those who practise patience, obedience,
and poverty climb to the third step; those who fight their vices and despise
the enemy triumph on the fourth with the Powers; faithful prelates who watch
for souls reign on the fifth with the Principalities; those who bring their
flesh under the spirit exult on the sixth with the Dominations; the
contemplatives, who offer God a peaceful dwelling, are companions of the
Thrones; the learned who pour back on God all they draw from him stand on
the eighth with the Cherubim. The ninth is love:
\begin{quote}
And they who love God with their whole mind and their whole heart, and who
place their whole being in that eternal fire, which is God, and who, being
made exceeding like unto Him, love Him no longer with their own, but with
His love \dots\ These shall come close to God, with nothing between Him and
them, on the ninth step, along with the Seraphim, between whom and God no
others are found. (ch.~16)
\end{quote}
The ladder follows the order of Gregory the Great, with the Principalities
above the Powers (\emph{Hom.\ in Evang.}\ 34.7), and it turns into a rule
of life the doctrine of Gregory and the Fathers that the elect are taken up
into the choirs of the angels (\S\ref{sec:fathers-after}).

\sectionguard\subsection{Joan of Arc and Saint Michael}\label{sec:sa-joan}

Joan of Arc gave an account of the archangel in her own words, under oath,
to the judges who condemned her. Her trial at Rouen in
1431 and the depositions taken for her rehabilitation were put into English
by T.\ Douglas Murray (1902). Her voices were Saint Catherine and Saint Margaret, but the first to come
was Michael. At the fourth public session, on 27 February, the judges pressed
her:
\begin{quote}
``What was the first Voice that came to you when you were about thirteen?''
``It was Saint Michael: I saw him before my eyes; he was not alone, but
quite surrounded by the Angels of Heaven. I came into France only by the
order of God.'' ``Did you see Saint Michael and these Angels bodily and in
reality?'' ``I saw them with my bodily eyes as well as I see you; when they
went from me, I wept. I should have liked to be taken away with them.''
\end{quote}
She would not describe him for their curiosity. Asked on 1 March whether he
came naked, she answered, ``Do you think God has not wherewithal to clothe
him?''; asked on 3 March whether Michael and Gabriel had human heads, ``I saw
them with my eyes; and I believe it was they as firmly as I believe there is
a God.'' On 15 March the examiners set the question of the discernment of
spirits before her, and her answers are those of a soul taught by the
archangel himself:
\begin{quote}
``How did you know it was Saint Michael?'' ``By the speech and language of
the Angels. I believe firmly that they were Angels.'' \dots\ ``When Saint
Michael came to me, he said to me: `Saint Catherine and Saint Margaret will
come to thee; follow their counsel; they have been chosen to guide thee and
counsel thee in all that thou hast to do: believe what they shall tell thee,
it is the order of Our Lord.'\,'' ``If the devil were to put himself in the
form or likeness of an angel, how would you know if it were a good or an
evil angel?'' ``I should know quite well if it were Saint Michael or a
counterfeit. The first time I was in great doubt if it were Saint Michael;
and I was much afraid. I had seen him many times before I knew it was Saint
Michael.''
\end{quote}
He had taught her, she said, ``so well, and it was so clear to me, that I
believed firmly it was he''; and his doctrine was simple: ``Above all things
he told me to be a good child, and that God would help me.'' On 17 March,
asked in what form he came, she said: ``In the form of a true honest man
[`prud homme']; of his dress and the rest I will say nothing more. As to the
Angels, I saw them with my eyes \dots\ I believe the deeds and words of Saint
Michael, who appeared to me, as firmly as I believe that Our Saviour Jesus
Christ suffered Death and Passion for us. And that which makes me believe
it, is the good counsel, comfort, and good doctrine which he has given me.''

Her test is the Church's test, the fruit of the spirit: a good angel gives
good counsel, comfort, and sound doctrine, and leads to God. She kept faith
with the archangel to the end. A priest present at the sermon of Saint-Ouen
deposed in 1452 that she prayed aloud ``to Saint Michael that he
would direct and counsel her,'' and that ``While they were tying her to the
stake she implored and specially invoked Saint Michael''; the physician
Guillaume Delachambre heard her at the last ``cry `Jesus' and \dots\ invoke
St.\ Michael; and then she perished in the flame.'' The archangel to whom,
as the Church sings, God has given the souls of the saints was the first
voice she heard and the last name she called.

\sectionguard\subsection{Frances of Rome and Her Archangel}\label{sec:sa-frances}

Frances of Rome, wife, mother, and foundress of the Oblates of Tor di
Specchi, lived the whole of her later life in the visible presence of an
angel. Her confessor, Giovanni Mattiotti, required her under obedience to
tell him her visions in detail, kept a daily record for some years, and
after her death wrote her life, which stands in the Bollandist collection;
Lady Georgiana Fullerton's life of 1855 is drawn from it and from the later
lives.

Her guardian first made himself known by correction. He was ``one day to
accompany her, not by an invisible presence only, as in the case of all
Christians, but, by a rare privilege of grace, in a visible form,'' and
before he showed himself he began to watch her conduct: ``At the least
imperfection in her conduct, before she had time to accuse and to condemn
herself, she felt the blow of a mysterious hand'' (Fullerton, ch.~3).
Tradition relates that after the death of her little son Evangelista he
appeared to her in her oratory with a companion brighter than himself, and
told her of the order of heaven: ``There are nine choirs of angels in
heaven, and the higher orders of angelic spirits instruct in the Divine
mysteries the less exalted intelligences.'' Of his companion he said, ``This
my companion is higher than I am in rank, as he is more bright and fair in
aspect. The Divine Majesty has assigned him to you as a guardian during the
remainder of your earthly pilgrimage'' (ch.~7). Evangelista departed; ``but
the archangel remained, and to the day of her death was ever present to her
sight.''

At her director's command she described him: ``His stature,'' she said,
``is that of a child of about nine years old; his aspect full of sweetness
and majesty; his eyes generally turned towards heaven \dots\ When I look upon
him, I understand the glory of the angelic nature, and the degraded
condition of our own'' (ch.~7). The light about him was so strong that ``At
night, and in the most profound darkness, she could always write and read
by the light of that supernatural brightness.'' In her combats with the
evil one he was her shield; the rays from his brow drove the demons off. The
angel was also her mirror and her rule. ``When she committed the slightest
fault, the angel seemed to disappear; and it was only after she had
carefully examined her conscience, discovered her failing, lamented and
humbly confessed it, that he returned''; he taught her to find the will of
God in the cares of her household, which she had been tempted to think lost
time, and he moderated her penances: ``Be not afraid, father; the archangel
will not allow me to proceed too far in that course'' (ch.~7). Her death was
his last errand. Asked by her confessor what she saw, she whispered, ``The
heavens open! The angels descend! The archangel has finished his task. He
stands before me. He beckons to me to follow him'' (ch.~14). In Frances is
made visible what the Catechism teaches of every believer: ``From infancy
to death human life is surrounded by their watchful care and intercession''
(CCC~336; \S\ref{sec:q113}).

\sectionguard\subsection{Ignatius of Loyola: The Good Angel and the Angel of
Light}\label{sec:sa-ignatius}

Ignatius gave the Church a school of prayer in which the angels, fallen and
unfallen, have a fixed place. His \emph{Spiritual Exercises} were put into
English from the Spanish autograph by Elder Mullan (1914).

The first exercise of the First Week begins with the first sin in the
world, the sin of the angels. The exercitant is to bring memory,
understanding, and will to bear on it:
\begin{quote}
I say to bring to memory the sin of the Angels, how they, being created in
grace, not wanting to help themselves with their liberty to reverence and
obey their Creator and Lord, coming to pride, were changed from grace to
malice, and hurled from Heaven to Hell. (First Week, first exercise, first
point)
\end{quote}
The whole doctrine of the fall is here in a sentence: created good and in
grace, free, fallen by pride, fixed in malice (\S\ref{sec:fall}; Lateran~IV,
DS~800). Its use is compunction: the exercitant sets ``the one sin of the
Angels'' beside his many sins and considers ``while they for one sin were
cast into Hell, how often I have deserved it for so many.'' The holy angels
appear at the end of the second exercise, in the cry of wonder at all
creatures that have borne with the sinner: ``the Angels, how, though they are
the sword of the Divine Justice, they have endured me, and guarded me, and
prayed for me'' (second exercise, fifth point). In the meditation on the Two
Standards the fallen angels are an army. The exercitant sees ``the chief of
all the enemy'' seated in the field of Babylon ``as in a great chair of fire
and smoke, in shape horrible and terrifying,'' and considers ``how he issues a
summons to innumerable demons and how he scatters them, some to one city and
others to another, and so through all the world, not omitting any
provinces, places, states, nor any persons in particular'' (Second Week,
fourth day). Against him stands Christ in a lowly place near Jerusalem,
sending out his friends.

Ignatius's rules for the discernment of spirits turn on the two angels who
move the soul. ``It is proper to God and to His Angels in their movements to
give true spiritual gladness and joy, taking away all sadness and
disturbance which the enemy brings on'' (Rules for the Second Week,
rule~1). God alone can console the soul without preceding cause; but ``With
cause, as well the good Angel as the bad can console the soul, for contrary
ends: the good Angel for the profit of the soul, that it may grow and rise
from good to better, and the evil Angel, for the contrary'' (rule~3). Then
comes the rule that makes Saint Paul's warning a science, ``Satan himself
transformeth himself into an angel of light'' (2 Cor 11:14):
\begin{quote}
It is proper to the evil Angel, who forms himself under the appearance of
an angel of light, to enter with the devout soul and go out with himself:
that is to say, to bring good and holy thoughts, conformable to such just
soul, and then little by little he aims at coming out drawing the soul to
his covert deceits and perverse intentions. (rule~4)
\end{quote}
The course of the thoughts reveals the spirit: if beginning, middle, and end
are all good, ``it is a sign of the good Angel''; if they end in something
bad or less good, or disquiet the soul and take away its peace, ``it is a
clear sign that it proceeds from the evil spirit, enemy of our profit and
eternal salvation'' (rule~5). The seventh rule gives the image: in those who go from good to better
``the good Angel touches such soul sweetly, lightly and gently, like a drop
of water which enters into a sponge; and the evil touches it sharply and
with noise and disquiet, as when the drop of water falls on the stone''
(rule~7). The Church's doctrine of the angels' action on the soul
(\S\ref{sec:q111}, \S\ref{sec:q114}) is here put to work in the soul's
choosing.

\sectionguard\subsection{Teresa of Jesus: The Angel of the
Transverberation}\label{sec:sa-teresa}

Teresa of Avila saw angels often, but seldom with the eyes. In the
twenty-ninth chapter of her \emph{Life} (trans.\ David Lewis, third
edition, 1904) she describes the grace that tradition calls her
transverberation, the piercing of her heart:
\begin{quote}
Our Lord was pleased that I should have at times a vision of this kind: I
saw an angel close by me, on my left side, in bodily form. This I am not
accustomed to see, unless very rarely. Though I have visions of angels
frequently, yet I see them only by an intellectual vision, such as I have
spoken of before. It was our Lord's will that in this vision I should see
the angel in this wise. He was not large, but small of stature, and most
beautiful---his face burning, as if he were one of the highest angels, who
seem to be all of fire: they must be those whom we call cherubim. Their
names they never tell me; but I see very well that there is in heaven so
great a difference between one angel and another, and between these and the
others, that I cannot explain it. (\emph{Life} 29.16)
\end{quote}
\begin{quote}
I saw in his hand a long spear of gold, and at the iron's point there seemed
to be a little fire. He appeared to me to be thrusting it at times into my
heart, and to pierce my very entrails; when he drew it out, he seemed to draw
them out also, and to leave me all on fire with a great love of God. The
pain was so great, that it made me moan; and yet so surpassing was the
sweetness of this excessive pain, that I could not wish to be rid of it. The
soul is satisfied now with nothing less than God. The pain is not bodily,
but spiritual; though the body has its share in it, even a large one.
(29.17)
\end{quote}
Her own manuscript names the angel's order as the cherubim; Lewis notes
that the earlier editors printed ``seraphim'' after a marginal suggestion of
her confessor Dominic Báñez, and restores her reading (29.16, n.~14). Either
way the angel is of the order that seems all fire, and the grace is of one
kind with the Seraph of La Verna: an angel of the highest choirs sent to
wound a soul with the love of God. Teresa's teaching on the angels is of a piece
with it. She insists that she sees the angels by intellectual vision and
that the distinctions among them are real and past telling. At another time,
rapt in the church, she saw the throne of God borne up and ``an exceedingly
great multitude of angels, who seemed to me more beautiful, beyond all
comparison, than those I had seen in heaven. I thought they were, perhaps,
the seraphim or cherubim, for they were very different in their glory, and
seemingly all on fire'' (39.32).

She knew the other side too. In the thirty-first chapter she recounts the
assaults of the devil and the weapon that served her best: ``I know by
frequent experience that there is nothing which puts the devils to flight
like holy water. They run away before the sign of the cross also, but they
return immediately: great, then, must be the power of holy water.'' The
reason is the Church's blessing: ``I have a joy in reflecting that the words
of the Church are so mighty, that they endow water with power'' (31.4).

\sectionguard\subsection{Francis de Sales and Peter Faber}\label{sec:sa-desales}

Francis de Sales brought friendship with the angels into the devout life of
the laity. The chapter of the \emph{Introduction to the Devout Life} on
honouring and invoking the saints (II.16; French edition of Lyon and Paris,
1832) begins with them:
\begin{quote}
\textit{Puisque c'est par le ministère des anges que nous recevons souvent
les bonnes inspirations de Dieu, c'est aussi par eux que nous devons lui
adresser nos aspirations.}
\end{quote}
Since it is by the ministry of the angels that we often receive God's good
inspirations, it is through them too that we should send him our
aspirations. After the saints and the Mother of God he gives his counsel:
\begin{quote}
\textit{Rendez-vous fort familière avec les anges: regardez-les comme
réellement présens à toutes vos actions, quoique d'une manière invisible.
Aimez surtout et respectez l'ange du diocèse où vous êtes, les anges des
personnes avec lesquelles vous vivez, et spécialement le vôtre: priez-les
souvent, offrez-leur de fréquentes louanges, et employez leur bon secours
dans toutes vos affaires, soit spirituelles, soit temporelles, afin qu'ils
coopèrent à vos intentions.} (II.16)
\end{quote}
Make yourself very familiar with the angels; look on them as really present
at all your actions, though invisibly. Love and respect above all the angel
of the diocese where you are, the angels of the persons with whom you live,
and especially your own; pray to them often, offer them frequent praise,
and use their good help in all your affairs, spiritual and temporal, so that
they may co-operate with your intentions. The counsel is a whole doctrine
of the angels' offices in small: the angels of dioceses and places, whom the
Fathers found in the angels of the nations (Dan 10:13; Deut 32:8,
Septuagint; \S\ref{sec:q113}), the guardians of our neighbours, and our own.

For an example he gives a saint of his own diocese. Peter Faber, first
priest and first companion of Ignatius, coming back from his labours in
Germany through the diocese of Geneva, where he was born, told how in heretical countries
\textit{il s'étoit toujours très-bien trouvé de saluer en arrivant dans une
paroisse les anges qui la protégeoient} --- he had always found it very
good, on arriving in a parish, to greet the angels who protected it, and
owed to that practice his escape from the snares of the heretics and the
docility of souls to the doctrine of salvation. A lady who heard him as a
girl repeated his words to Francis more than sixty years later, and the
bishop consecrated an altar in the village where Faber was born (II.16).
The first Jesuit greeting the angels of each parish and the bishop of
Geneva counselling the devout to love the angel of their diocese stand in
one line: the guardian angels of places are the Church's allies in the
work of the Gospel.

\sectionguard\subsection{Alphonsus Liguori}\label{sec:sa-alphonsus}

Alphonsus places the angels above all at the hour of death. In his sermon
for the Sixth Sunday after Epiphany, on the death of the just, he meets the fear of the soul that dreads the last assaults of
hell. God does not abandon his servants then; when Eliseus's servant saw
the city ringed with enemies, the prophet showed him the angels sent to
defend it (4~Kings 6:16--17). So now:
\begin{quote}
The powers of hell will assail the dying Christian; but his angel guardian
will come to console him. His patrons, and St.\ Michael, who has been
appointed by God to defend his faithful servants in their last combat with
the devils, will come to his aid. (\emph{Sermons for all the Sundays},
s.~11, trans.\ Grimm, p.~125)
\end{quote}
In his sermon on death for the Fifteenth Sunday after Pentecost, the
priest at the bedside bids the dying man say after him, ``St.\ Joseph, St.\
Michael, the archangel, my holy angel-guardian, and all ye saints in
paradise, pray to God for me'' (s.~44, p.~458). In the \emph{Glories of Mary}
the archangel serves his Queen at the same hour. Alphonsus quotes, under
Bonaventure's name, the \emph{Speculum Beatae Virginis}: Mary ``sends
without delay the prince of the heavenly court, St.\ Michael, with all the
angels, to defend her dying servants against the temptations of the devils,
and to receive the souls of all who in a special manner and perseveringly
have recommended themselves to her'': \latin{Michael, dux et princeps militiæ
cœlestis, cum omnibus spiritibus administratoriis, tuis, Virgo, paret
præceptis} (\emph{Glories of Mary} I.2.3). The office that made Francis love
Michael, the presenting of souls, is here set within the Queenship of Mary
(\S\ref{sec:christ-mary}).

\sectionguard\subsection{John Vianney, the Curé of Ars}\label{sec:sa-vianney}

The Curé of Ars spoke of the angels as he spoke of everything, in plain
words to his parishioners; Alfred Monnin gathered his catechisms and sayings
(\emph{The Spirit of the Curé of Ars}, English ed.\ 1865). The guardian
angel is always by us: ``When we are going along the streets, let us fix our eyes on our Lord carrying His cross before us; on
the Blessed Virgin, who is looking at us; on our guardian angel, who is by
our side'' (Catechism on the Cardinal Virtues). He keeps our account: ``We
always have two secretaries: the devil, who writes down our bad actions, to
accuse us of them; and our good angel, who writes down our good ones, to
justify us at the Day of Judgment'' (pt~I, introduction). He rejoices in
the pure: ``What joy is it to the guardian angel to conduct a pure soul!'' (Catechism
on the Prerogatives of the Pure Soul); and ``How happy is a guardian angel
who leads a beautiful soul to the holy table!'' (Catechism on Frequent Communion). His reverence
for the priesthood set even the angels second: ``If I were to meet a priest
and an angel, I should salute the priest before I saluted the angel. The
latter is the friend of God; but the priest holds His place'' (Catechism on
the Priesthood).

His biographers record that for many years the devil harassed him by night
with noise and violence. Kathleen O'Meara's life gives his own account of
the first night: he went back to bed ``commending myself to God, to the
Blessed Virgin, and my Guardian Angel,'' and when he knew the cause he
prayed God ``to be near me with His angels when the enemy came to torment
me.'' He met the trial with the Sign of the Cross and with laughter: ``The
devil is cunning \dots\ but he is not strong; one Sign of the Cross puts him
to flight''; and of his tormentor, whom he called only \emph{le grappin},
``He is more powerful than the grappin'' --- meaning God (O'Meara, \emph{The
Curé of Ars}, ch.~9). O'Meara sets him beside Anthony and Hilarion, whom
the demons tormented in the desert (\S\ref{sec:assaults}); the lesson the
Curé drew is the Church's own, that the enemy can frighten and cannot harm
those whom God keeps.

\sectionguard\subsection{Gemma Galgani and Her Guardian Angel}\label{sec:sa-gemma}

Gemma Galgani of Lucca lived with her guardian angel as a child with an
elder brother. Her director, the Passionist Father Germanus, devotes a chapter of her
life to him (English trans.\ A.\ M.\ O'Sullivan, 1913). He records that she saw the angel, touched him, and
spoke with him ``as one friend would with another.'' ``Jesus,'' she said,
``has not left me alone; He makes my guardian angel stay with me always'';
and the angel told her, ``Yes I shall be thy guide and inseparable
companion. Do you not know who it is that gave me charge of thee? It is the
merciful Jesus'' (Germanus, \emph{Life}, ch.~20). They prayed together, the
angel sometimes ``raised in the air with outspread wings, and his hands
extended over her or else joined in attitude of prayer,'' sometimes kneeling
beside her; they said the psalms in alternation and vied with each other in
aspirations; at meditation the angel opened to her the Passion. Germanus writes
as an eyewitness: ``I myself have many times assisted at these meditations of
Gemma with her angel,'' and whenever she raised her eyes to the angel she
lost the use of her senses.

The angel was her master in perfection. Once he made her write at his
dictation, ``like a child at school'':
\begin{quote}
``Remember child, that one who truly loves Jesus, speaks little and endures
much. I command thee on the part of Jesus never to give thy opinion unless
it is asked; never to maintain thy opinion, but be silent at once \dots\
Remember to guard thine eyes and reflect that the mortified eye shall behold
the beauties of heaven.'' (ch.~20)
\end{quote}
He corrected her with severity: ``My Angel is a little severe; but I am glad
of it. During the past days he brought me to order as often as three or four
times a day.'' He carried her messages and brought her home from church when
she had forgotten the hour; when the devil beat her at her evening prayer,
the angel lifted her to her bed and watched by her pillow. Her love for him
was simple: ``Dear Angel,'' she would say, ``I so love you!'' ``And why?''
``Because you teach me to be good, and to keep humble, and to please
Jesus.'' Her director tested her by forbidding familiarity, and she obeyed
at once: ``From this day forward, I will always use the right terms when
speaking to the Angel, and I will show him every reverence.'' Obedience was
the proof the director sought and the angel himself required; when she
feared to disobey by listening to him he answered, ``No, because I am sent by
Him Who has power to send me.'' Once he turned her from creatures to God: ``It is my wish that thy conversation be no longer with creatures; when
you do wish to speak, speak with Jesus, and with thy Angel'' (ch.~20).
Germanus finds the parallel in Scripture, in Raphael and the young Tobias
(Tob 12:12--15), and the friendship of Gemma with her angel shows in visible
form what every guardian does unseen.

\sectionguard\subsection{Thérèse of the Child Jesus}\label{sec:sa-therese}

Thérèse of Lisieux, who died on 30 September 1897, the day after the feast
of Saint Michael, wrote poems in which the angels take their part. In February 1897, seven
months before her death, she addressed her guardian angel. The poem is printed in \emph{Histoire
d'une âme} (edition of 1912, pp.~427--428):
\begin{quote}
\textit{Glorieux gardien de mon âme,\\
Toi qui brilles dans le beau ciel\\
Comme une douce et pure flamme,\\
Près du trône de l'Éternel ;\\
Tu viens pour moi sur cette terre,\\
Et m'éclairant de ta splendeur,\\
Bel Ange, tu deviens mon frère,\\
Mon ami, mon consolateur !}
\end{quote}
Glorious guardian of my soul, you who shine in the fair heaven like a sweet
and pure flame near the throne of the Eternal, you come to this earth for
me, and lighting me with your splendour, fair Angel, you become my brother,
my friend, my consoler. The second stanza is the little way in the
angel's hands: \textit{Connaissant ma grande faiblesse, / Tu me diriges par
la main ; / Et je te vois, avec tendresse, / Ôter la pierre du chemin} ---
knowing my great weakness, you lead me by the hand, and I see you tenderly
take the stone out of the road (Ps 90[91]:12). She sends him as her
messenger, since he crosses space \textit{plus promptement que les éclairs},
more swiftly than the lightning, to dry the tears of those she loves, and
asks his zeal for souls: \textit{Je veux, pendant ma courte vie, / Sauver mes
frères les pécheurs ; / Ô bel Ange de la patrie, / Donne-moi tes saintes
ardeurs} --- I wish, in my short life, to save my brothers the sinners; fair
Angel of the fatherland, give me your holy ardour. The last stanzas divide
the goods of heaven and earth between them:
\begin{quote}
\textit{À toi, le royaume et la gloire,\\
Les richesses du Roi des rois.\\
À moi, le Pain du saint ciboire,\\
À moi, le trésor de la Croix.}
\end{quote}
Yours the kingdom and the glory, the riches of the King of kings; mine the
Bread of the holy ciborium, mine the treasure of the Cross. The angel has
the glory; the child on earth has what the angel cannot have, the Eucharist
and the Cross, and she offers her poverty to the Trinity through him.

In her piece on the angels at the crib, the angel of the Child
Jesus gazes on the Word made man in the manger; in Susan Emery's translation
of 1907 he sings, ``O Child, whose light doth blind the sight / Of angels in
high heaven divine! / Thou'rt come to save the world to-night, / And who can
fathom that love of Thine?'' and promises, ``Thee will I guard by day and
night'' (\emph{Poems of Sr.\ Teresa}, ``The Angels of the Crib''). The
guardian of the Child in that piece and her own guardian in her poem are the
same service, the angels' care for the little ones whose angels see the
Father's face (Matt 18:10).

\sectionguard\subsection{John Henry Newman}\label{sec:sa-newman}

Newman preached on the angels on the feast of Saint Michael and All Angels
while he was still a clergyman of the Church of England, and as a priest
of the Oratory wrote the \emph{Dream of Gerontius}. The sermon, ``The Powers of Nature'' (\emph{Parochial and Plain Sermons} II.29),
takes for its text ``Who maketh His Angels spirits, his ministers a flaming
fire'' (Ps 103[104]:4) and sets it against the temper of the age:
\begin{quote}
There have been ages of the world, in which men have thought too much of
Angels, and paid them excessive honour; honoured them so perversely as to
forget the supreme worship due to Almighty God. This is the sin of a dark
age. But the sin of what is called an educated age, such as our own, is just
the reverse: to account slightly of them, or not at all; to ascribe all we
see around us, not to their agency, but to certain assumed laws of nature.
(pp.~358--359)
\end{quote}
Newman answers from Scripture with the doctrine that the angels move the
visible creation (\S\ref{sec:q110}). Reason sees that spirit moves man and beast, but
``how do the wind and water, earth and fire, move? Now here Scripture
interposes, and seems to tell us, that all this wonderful harmony is the
work of Angels.'' An angel gave the pool of Bethesda its healing, angels
wrought the fires of Sinai, restrained the four winds, and smote the camp of
Sennacherib. ``Nature is not inanimate; its daily toil is intelligent; its
works are duties'' (pp.~360--361). Then:
``Every breath of air and ray of light and heat, every beautiful prospect,
is, as it were, the skirts of their garments, the waving of the robes of
those whose faces see God in heaven'' (p.~362). The science of nature is
useful ``in subordination to that higher view''; but the man who takes the
order he discerns for ``fixed laws, self-caused and self-sustained'' forgets
``the thousands and ten thousands of His unseen Servants.'' The doctrine
humbles, consoles, and draws the soul to heaven: ``it is a great comfort to
reflect that, wherever we go, we have those about us, who are ministering to
all the heirs of salvation, though we see them not'' (Heb 1:14); ``The very
lowest of His Angels is indefinitely above us in this our present state; how
high then must be the Lord of Angels!''; and ``if we attain to heaven, we
shall become the fellows of the blessed Angels'' (pp.~365--366).

\subsubsection{The Dream of Gerontius}

The \emph{Dream of Gerontius} (1865) follows a soul from its death-bed to
Purgatory in the arms of its guardian angel. When the soul wakes after
death it hears the angel singing:
\begin{quote}
My work is done,\\
My task is o'er,\\
And so I come,\\
Taking it home,\\
For the crown is won,\\
Alleluia,\\
For evermore. (\S~2)
\end{quote}
``My Father gave / In charge to me / This child of earth / E'en from its
birth, / To serve and save.'' The soul recognizes in him ``a member of that
family / Of wondrous beings, who, ere the worlds were made, / Millions of
ages back, have stood around / The throne of God:---he never has known
sin.'' The angel himself states the privilege of the guardians: ``More than the Seraph
in his height of place, / The Angel-guardian knows and loves the ransom'd
race'' (\S~2). On the way to judgment they pass the choirs of Angelicals,
``Least and most childlike of the Sons of God,'' who sing ``Praise to the
Holiest in the height, / And in the depth be praise'' and tell of God's
``elder race'' made ``To battle and to win, / Without the chastisement of
pain, / Without the soil of sin'' (\S~5). Before the throne the soul hears
the prayers of its friends saying the \latin{Subvenite} on earth, and the
guardian points to ``the great Angel of the Agony, / The same who
strengthen'd Him, what time He knelt / Lone in that garden shade, bedew'd
with blood'' (\S~6; Luke 22:43), who pleads for the souls in Purgatory. At
the last the guardian lowers the soul into the cleansing waters with a
promise:
\begin{quote}
Softly and gently, dearly-ransom'd soul,\\
In my most loving arms I now enfold thee,\\
And, o'er the penal waters, as they roll,\\
I poise thee, and I lower thee, and hold thee. \dots\\
Farewell, but not for ever! brother dear,\\
Be brave and patient on thy bed of sorrow;\\
Swiftly shall pass thy night of trial here,\\
And I will come and wake thee on the morrow. (\S~7)
\end{quote}
The guardian's office does not end at death; he brings the soul to
judgment and will bring it out of Purgatory, as the Church prays in the
\latin{In paradisum} (CCC~335; \S\ref{sec:liturgy}).

\sectionguard\subsection{The Saints' Witness and Its Grade}\label{sec:sa-table}

\begin{fourstable}{Saint}{Witness and locus}{Devotion or teaching}{Grade}
Francis of Assisi & Bonaventure, \emph{Legenda maior} prol.~1--2; II.8; IX.3; XIII.1--5 & Lent of Saint Michael; love of Michael for his office of presenting souls; the Seraph of La Verna and the stigmata & Pious tradition; the angels' ministry to the saved, Church teaching (CCC~336)\\
Bonaventure & \emph{Itinerarium} prol.~2--3; \emph{Sermones de sanctis}, De sanctis Angelis 1, 4, 5 & The six wings as six illuminations; nine orders in three hierarchies with their offices; Raphael, Gabriel, Michael as remedies; angels as general guardians of souls & Common teaching on the hierarchies (Dionysius, \S\ref{sec:ch}); Church teaching on the guardians (CCC~336)\\
Thomas Aquinas & Tocco, \emph{Vita} c.~10; Calo, \emph{Vita} 7; Pius XI, \emph{Studiorum Ducem} (1923) & The angels gird him with chastity; the Angelic Doctor; the Angelic Warfare; the prayer \latin{Creator ineffabilis} & Pious tradition for the girding; the title Angelic confirmed by Pius V (\emph{Studiorum Ducem})\\
Gertrude & \emph{Insinuationes}, Kenmare trans.\ pt~III ch.~21; pt~IV chs.~2, 54 & Guardian presents the soul at Mass; the choirs at the Divine Office; the Eucharist offered in honour of the angels & Pious tradition\\
Mechtild & \emph{Select Revelations} ch.~16 & Nine steps of the orders joined to nine ways of life; the Seraphim nearest God & Pious tradition; order of Gregory (\S\ref{sec:fathers-after})\\
Joan of Arc & Trial of 1431, 27 Feb., 1, 3, 15, 17 Mar.; rehabilitation depositions (Murray, 1902) & Saint Michael, first of her voices; his good counsel as the sign of a good angel; invoked at her death & Pious tradition\\
Frances of Rome & Fullerton, \emph{Life} chs.~3, 7, 14 (after Mattiotti) & A visible archangel as guardian; nine choirs and illumination of the lower by the higher & Pious tradition; the angels' care from infancy to death Church teaching (CCC~336)\\
Ignatius of Loyola & \emph{Spiritual Exercises}, first exercise; Two Standards; Rules for the Second Week 1--7 & The sin of the angels; the demons sent through the world; good and evil angel in the soul; the angel of light & Defined that the demons fell of themselves (DS~800); common teaching on discernment (2 Cor 11:14)\\
Teresa of Jesus & \emph{Life} 29.16--17; 31.4; 39.32 & The cherub of the transverberation; intellectual vision of angels; holy water puts the devils to flight & Pious tradition\\
Francis de Sales & \emph{Introduction} II.16 & Familiarity with the angels; the angels of the diocese, of neighbours, and one's own; Faber's greeting of the angels of parishes & Common teaching (Dan 10:13; \S\ref{sec:q113})\\
Alphonsus Liguori & \emph{Sermons} 11, 44; \emph{Glories of Mary} I.2.3 & Guardian and Michael assist the dying; Michael obeys the Queen of heaven & Common teaching; Pious tradition on the \emph{Speculum}\\
John Vianney & Monnin, \emph{Spirit} (1865); O'Meara, ch.~9 & The guardian angel at our side and at Communion; the priest saluted before the angel; the devil's harassment & Pious tradition\\
Gemma Galgani & Germanus, \emph{Life}, ch.~20 & Visible guardian as companion, master, and messenger, under obedience & Pious tradition\\
Thérèse & \emph{Histoire d'une âme} (1912), pp.~427--428; ``The Angels of the Crib'' (Emery) & The guardian as brother, friend, consoler; the angels' glory and the Cross and the Host & Pious tradition\\
John Henry Newman & \emph{Parochial and Plain Sermons} II.29; \emph{Dream of Gerontius} \S\S~2, 5--7 & The angels move the visible creation; the guardian brings the soul to judgment and Purgatory & Common teaching on the angels' government of bodies (\S\ref{sec:q110}); Pious tradition in the poem\\
\end{fourstable}
